\section{Escape Dynamics}
\label{sec:Escape}

This section formalizes the escape criterion for the quadratic family.

\subsection{Escape Radius}

We define a bound $R(c) = \max(2, |c|+1)$.
The main lemma (\texttt{escape\_lemma}) states that if $|z| > R(c)$, then $|f_c(z)| > |z|$ and the orbit escapes to infinity geometrically.

\[ |f_c(z)| = |z^2 + c| \ge |z|^2 - |c| = |z|(|z| - \frac{|c|}{|z|}) \]
Since $|z| > 2$ and $|z| > |c|$, the factor is greater than 1.

\subsection{Formalization}

\textbf{Lemma} (\texttt{norm\_orbit\_ge\_of\_norm\_ge\_R}): Points starting outside the escape radius remain outside and their norm grows.

\begin{lstlisting}[language=Lean]
/-- Points outside the escape radius escape to infinity. -/
lemma norm_orbit_ge_of_norm_ge_R (c z : $\mathbb{C}$) (n : $\mathbb{N}$) (h : $\|$z$\|$ > R c) :
    $\|$orbit c z n$\|$ $\ge$ $\|$z$\|$
\end{lstlisting}

This property is crucial for defining the Green's function, as it ensures that the series defining the potential converges for points outside $K_c$.
